\documentclass[11pt,a4paper]{article}
\usepackage[margin=25mm]{geometry}
\usepackage{amsmath,amssymb,amsthm,hyperref}
\newtheorem{theorem}{Theorem}
\newtheorem{lemma}[theorem]{Lemma}
\newcommand{\gen}[1]{\langle#1\rangle}
\title{An endpoint anchor in the imprimitive rank bound}
\author{Human Compression project: bounded L044 proof component}
\date{}
\begin{document}\maketitle
We distinguish the literal general assertions in L044
Propositions~5.12(c) and~5.13 from the narrower assertion needed
for a supported interval which extends in a larger string action.
All permutations below act on specified points; parity is measured
in that action. An \emph{endpoint anchor} for a transitive marking
$(s_0,\ldots,s_{d-1})$ is a point fixed by all generators except
one endpoint generator. Transitivity implies that endpoint moves it.

\begin{theorem}[A uniform obstruction to the unanchored assertion]
For every odd $b\geq3$, there is an even transitive imprimitive
string C-group on $2b$ points of rank $b$, with connected
Schl\"afli diagram $[4,3,\ldots,3]$, and every maximal omission
intransitive. Neither endpoint has an anchor.
\end{theorem}
\begin{proof}
Write the points as $(j,\epsilon)$ with $1\leq j\leq b$ and
$\epsilon\in\mathbb F_2$. In the signed permutation group
$V\rtimes S_b$, $V=\mathbb F_2^b$, let $e_j$ flip the $j$th
pair and let $u=(1,\ldots,1)$. Define the $S_b$-equivariant map
\[
 f(v)=v+\Big(\sum_jv_j\Big)u.
\]
As $b$ is odd, $\ker f=\gen u$ and $\operatorname{im}f=E$,
the even-weight subspace. Take
\[
 s_0=f(e_1)=u+e_1,\qquad
 s_j=(j\ j+1)(j'\ (j+1)')\quad(1\leq j<b).
\]
Here primes distinguish the two points in a pair.
These generate $E\rtimes S_b$, of order $2^{b-1}b!$.
Every generator is even. The diagonal $S_b$ is transitive on
pair labels, and $E$ has a vector flipping any specified pair,
so the point action is transitive. The pairs are a nontrivial
block system. The $s_j$ for $j\geq1$ are the ordinary adjacent
transpositions on pair labels. Further, $s_0$ commutes with
$s_j$ for $j\geq2$, while $(s_0s_1)^2=e_1+e_2$ has order two.
This gives the displayed connected diagram.

Here is a full intersection-property proof, independent of a
polytope quotient criterion. For a subset $I\subseteq\{0,\ldots,b-1\}$,
let $t(I)$ be the largest $t$ with $\{0,\ldots,t-1\}\subseteq I$,
put $V_t=\operatorname{span}(e_1,\ldots,e_t)$, and let
$P_I=\gen{(j\ j+1):j\in I,\ j\geq1}$.
Direct conjugation of $f(e_1)$ by these permutations gives
\[
                     G_I=f(V_{t(I)})\rtimes P_I.
\]
The prefix spaces $f(V_t)$ are nested, hence their pairwise
intersections are the smaller prefix space. The groups $P_I$
are products of symmetric groups on intervals of pair labels;
intersecting two such interval partitions gives
$P_I\cap P_J=P_{I\cap J}$. Finally
$t(I\cap J)=\min(t(I),t(J))$. Uniqueness of vector and permutation
coordinates now proves $G_I\cap G_J=G_{I\cap J}$ for every
$I,J$, including the full set. The marking is therefore a
string C-marking.

Omitting $s_0$ leaves the two natural $S_b$ point orbits.
Omitting $s_j$ for $j\geq1$ preserves the nonempty disjoint
sets of pairs $\{1,\ldots,j\}$ and $\{j+1,\ldots,b\}$.
Thus all maximal omissions are intransitive. The first endpoint
omission has no fixed point. The last contains both $f(e_1)$
and $f(e_2)$; their supports together meet every pair, so it
also has no fixed point. There is no endpoint anchor.
\end{proof}

For odd $b\geq7$, this rank $b$ exceeds
$\lfloor(2b-1)/2\rfloor=b-1$. Thus it contradicts the literal
unanchored conclusion of source Proposition~5.13 (take its
$m=2b$ and any ambient $n>2b$), and the
$(a,d)=(2,b)$ clause of Proposition~5.12(c). It does not
contradict the source's main assertion about $A_n$: the group
here is imprimitive and proper. It also lies outside our
accepted unique-crossing Proposition~5.7 setting. Indeed,
$s_0$ has $b-1$ eligible edges between the two omission orbits,
and each $s_j$, $j\geq1$, has both of its two edges crossing
the invariant pair-label partition. No colour has just one
eligible edge. This is the odd hemicube family already appearing in the source
itself, Corollary~4.19(b), physical pages~21--22. The explicit
marked intersection proof above verifies its relevant properties;
no new-family, historical-novelty or question-resolution claim is made.

\begin{lemma}[The anchor supplied by an extension]
Let $H=\gen{s_0,\ldots,s_{d-1}}$ be supported on $X$, and let
an involution $t$ centralize $s_0,\ldots,s_{d-2}$. If $t$ takes
some $c\in X$ to a point outside $X$, then $c$ is fixed by
$s_0,\ldots,s_{d-2}$. If $H$ is transitive on $X$, $c$ is an
anchor at its last endpoint.
\end{lemma}
\begin{proof}
For $j\leq d-2$, commutation and the fact that $s_j$ fixes
$c^t\notin X$ give $(c^{s_j})^t=(c^t)^{s_j}=c^t$, hence
$c^{s_j}=c$. If the last generator also fixed $c$, all of
$H$ would fix it, contrary to transitivity on a nontrivial $X$.
\end{proof}

\begin{theorem}[The anchored imprimitive bound]
Let $H\leq A_N$, $N\geq12$, be a transitive imprimitive
string C-group with connected noncommuting diagram, all maximal
omissions intransitive, and an endpoint anchor. Then
\[
                         d\leq\left\lfloor\frac{N-2}{2}\right\rfloor.
\]
\end{theorem}
\begin{proof}
Choose a block system with block size $a\geq2$ and block
number $b\geq2$. Let $M$ be a minimal distinguished subfamily
generating the complete block action, and let $c$ be its actual
number of point orbits. The accepted actual-orbit forest
argument, now on one transitive orbit, gives
\[
       d-|M|\leq c-1,\qquad c\leq a,\qquad |M|\leq b-1.
\]
For clarity, the forest selects one edge of each omitted
colour between different maximal-omission orbits and contracts
actual $\gen M$-orbits; any cycle would give a path avoiding
one selected colour. The last inequality is Whiston's bound,
with no quotient IP assumption. Consequently $d\leq a+b-2$.
If $a,b\geq3$, then $N\geq12$ excludes $(a,b)=(3,3)$, and
$(a-2)(b-2)\geq2$ gives the claimed bound.
It remains to exclude $d=a$ when $b=2$, and $d=b$ when $a=2$.

Suppose first $b=2$. Reverse the marking so that an anchor $p$
is moved only by $s_0$. Every other generator fixes $p$, hence
preserves the two blocks. Transitivity forces $s_0$ to exchange
them; denote it by $\tau$. It is a fixed-point-free product of
$a$ transpositions, so evenness forces $a$ even, with $a\geq6$.
If $d=a$, the $a-2$ generators after $x=s_1$ centralize $\tau$
and fix $p$. They act faithfully and diagonally on the other
$a-1$ row pairs, and their independent length is maximal.
Thus they generate the full $S_{a-1}$ by Whiston.
The accepted maximal independent involutory-string theorem
makes this the ordinary adjacent-transposition path:
$a-1\geq5$ is odd, so the outer degree-six exception is absent.
Its near-end omission is the natural $S_{a-2}$ on row labels,
fixing the anchor row and one unused row $q$.
Since $x$ commutes with that omission, preserves the two
blocks and fixes $p$, it is trivial on the block containing
$p$. On the opposite block it is either trivial or
$(p'\ q')$. Transitivity requires the latter: every other
generator except $\tau$ fixes $p'$ if $x$ does.
Thus $x$ would be odd, a contradiction.

Suppose instead $a=2$, so $b\geq6$, and assume $d=b$.
All three forest inequalities must be equalities. The block
image is $S_b$, and the two actual $\gen M$-orbits are natural
copies meeting each block in one point. Hence $\gen M$ is
faithfully the diagonal $S_b$. Its marking inherits IP.
It is connected: otherwise IP and string commutation would
split $S_b$ as a nontrivial direct product. Therefore the
$b-1$ members of $M$ are consecutive, and the single remaining
generator $x$ is at an end of the full marking.

Use the complete accepted maximal symmetric-string classification.
For $b\geq7$ the marking of $M$ is ordinary; for $b=6$ it
may also be the outer image of an ordinary marking.
The image of $x$ on the blocks centralizes the near-end
omission of $M$. In the ordinary case that omission is
$S_{b-1}$ fixing one block, whose centralizer in $S_b$ is
trivial. In the outer degree-six case it is the transitive
$S_5$ on six blocks, with trivial centralizer by the accepted
endpoint interface. Thus $x$ has trivial block image.
It is a vector of flips of the two-point blocks.

In the outer case, commutation with the transitive omission
forces that vector to be constant. It then centralizes all
of $M$, contrary to the connected diagram. In the ordinary
case the vector is constant on the $b-1$ blocks of the omission
and independent on its remaining block. Connectedness leaves
exactly the vectors $e_1$ and $u+e_1$, after row relabelling.
The first is odd. The second is even precisely when $b$ is
odd, and is exactly the marking of Theorem~1, up to reversal.
Neither endpoint of that marking has an anchor. This is the
last contradiction.
\end{proof}

\paragraph{Source and paid boundaries.}
The source is Cameron--Fernandes--Leemans--Mixer,
\emph{Highest rank of a polytope for $A_n$},
arXiv:1605.09173v1, physical pages~27--28. The basic forest
bound recovers Proposition~5.12(a,b); its part~(c) needs the
explicit correction above. Theorem~3 replaces the imprimitive
step only where its anchor hypothesis is actually proved.
Lemma~2 supplies that hypothesis for a supported interval that
extends to an external point; it is not an automatic assertion
for a quotient, a sesqui-extension, or an arbitrary smaller group.
The primitive and alternating branches of a full replacement
for Proposition~5.13 require their separate hypotheses and
accepted proofs. We retain the complete L047 common extremal
proofs, its bounded $S_5/S_6$ certificate and the actual-orbit
forest proof as paid dependencies. No finite experiment is a
premise of these new arguments. No full L044 recovery, new
formal declaration, new original credit or current-open-question
claim follows from this bounded component.
\end{document}
